\section{八、把片上网络落实为分包、缓冲与前进条件}

前面的 Crossbar 与 NoC 说明了拓扑，这里进一步追踪一笔事务在路由器里的实际状态。设链路宽度为 128 bit，也就是每个 flit 搬运 16 B。一笔携带 64 B Cache line 的写请求若使用两个头 flit 描述地址、源节点、事务 ID、长度与属性，再使用四个数据 flit，就形成六个 flit：head 建立路由，四个 body 搬数据，tail 释放本包占用的路由和虚通道状态。

\topic{wormhole 转发把缓冲需求变小，也把阻塞传得更远}

store-and-forward 路由器要等整包到齐后才转发，六个 flit 至少占六个槽；wormhole 路由让 head flit 选定下一跳后，后续 flit 可以流水跟进，每个路由器只需保存少量 flit。代价是 tail 尚未离开前，本包占用的通道仍不能交给冲突流量；若下游停顿，body 会像一列火车一样跨多个路由器排队，背压最终传回源节点。

无竞争时，六个 flit 穿过三跳、每跳路由与链路合计一拍，可以用简化模型估算：head 到目标需 3 拍，tail 比 head 晚发 5 拍，因此末 flit 在第 $3+5=8$ 拍到达。这个 8 拍只含网络流水，不含源端排队、目标处理和返回事务。实际时延应拆成
\[
T_{\text{end-to-end}}
=T_{\text{source queue}}+T_{\text{serialization}}
+T_{\text{hops}}+T_{\text{contention}}+T_{\text{target}} .
\]
只报告平均 hop 数会漏掉最容易形成长尾的源端排队和热点竞争。

\topic{credit 是一份必须守恒的远端缓冲账本}

credit 流控中，发送端的计数表示远端虚通道当前承诺的空槽数。发送一个 flit 就减一，远端真正释放槽位后才返还一。若远端有 8 个槽、credit 往返延迟为 6 拍，而链路每拍最多发送一个 flit，8 个 credit 足以覆盖流水；若只有 4 个，持续流量的长期上限约为 $4/6$ flit/拍，即使物理链路能做到 1 flit/拍也会周期性空闲。

\begin{longtable}{L{0.17\linewidth}L{0.25\linewidth}L{0.28\linewidth}L{0.20\linewidth}}
\toprule
事件 & 本地 credit & 远端槽位 & 不变量 \\
\midrule
初始化 & 8 & 8 个空槽 & credit 不超过已承诺容量 \\\tablerowrule
连续发送 4 flit & 4 & 4 个空槽 & 每次发送恰好消耗一个 credit \\\tablerowrule
远端消费 2 flit & 4 & 6 个空槽 & 返还消息尚在途中 \\\tablerowrule
返还到达 & 6 & 6 个空槽 & 本地重新获得发送许可 \\
\bottomrule
\end{longtable}

计数器位宽、复位初值和链路重训练都必须一致。若发送端在远端尚未清空旧 flit 时把 credit 直接恢复到满值，就会把同一个槽位计算两次；若返还消息丢失而没有重同步机制，容量虽仍存在，链路却会永久停在零 credit。

\topic{虚通道解决的是资源依赖，不是增加物理带宽}

虚通道让一条物理链路保存多组独立队列和 credit。它不能让导线一拍传两个 flit，却能防止一个被阻塞的包占住所有缓冲。最重要的用途是切断通道依赖环。例如请求包占用通道 R 并等待响应包，而响应包又必须经过被请求流量填满的 R，便可能形成 R 等 R 的环。把请求放在虚拟网络 VN0、响应放在 VN1，并规定响应从不再申请请求通道，就能让依赖图保持无环。

死锁、活锁和饥饿要分开：死锁是资源形成环，参与者都无法前进；活锁是包持续移动却始终到不了目标；饥饿是某个请求长期输掉仲裁。增加年龄优先可以缓解饥饿，却不能拆除通道依赖环；增加超时可以发现死锁，却不能证明协议不会死锁。

\section{九、顺序、QoS 与端到端完成语义}

NoC 可以让不同路径、不同虚通道和不同目标的包乱序到达。端点因此要明确顺序域：同一事务的 flit 不能错序；同一 ID 的请求是否保持顺序由协议规定；不同 ID 若允许重排，发起者必须按 ID 匹配响应。屏障事务还要等此前所属顺序域内的访问达到规定观察点，不能只因为屏障包先走一条空闲路径就提前完成。

QoS 也不能只看仲裁权重。设 1 GHz NoC 中一笔六 flit 包无竞争需 8 ns，目标服务需 20 ns，正常总时延约 28 ns。若热点端口前方已有五个同样的包且不可抢占，额外排队可能接近 $5\times6=30$ ns；若系统给实时流量规定 80 ns 上限，就必须同时约束注入率、最长包、仲裁等待和目标服务时间。权重只控制长期份额，不自动给出单笔最坏时延。

\begin{keyidea}
互连正确性最终由四组不变量组成：缓冲不能溢出或重复记账，通道依赖不能成环，响应必须能归属原请求，顺序与屏障必须在端点观察点闭合。拓扑图只能说明“可能往哪里走”，这些状态才说明“为何一定能正确到达”。
\end{keyidea}

\begin{chapterexercise}{credit 吞吐}
远端虚通道有 5 个槽，credit 往返延迟为 8 拍，链路峰值为每拍 1 flit。估算只受 credit 限制时的长期吞吐上限，并说明怎样达到满速。
\exercisecheck{检查要点}{窗口至少覆盖带宽时延积}
\end{chapterexercise}

\begin{chapterexercise}{通道依赖}
请求通道等待响应，响应又只能申请被请求占用的同一缓冲。画出依赖环，并给出虚拟网络划分规则。
\exercisecheck{检查要点}{优先级不能替代无环资源依赖}
\end{chapterexercise}

\begin{chapteranswer}{credit 吞吐}
上限约为 $5/8=0.625$ flit/拍。要达到每拍 1 flit，credit 窗口至少覆盖 8 拍往返，即至少 8 个可用槽，或缩短返还延迟。
\answercheck{必须保留的检查点}{credit 是远端容量承诺，不是任意加大的性能旋钮}
\end{chapteranswer}

\begin{chapteranswer}{通道依赖}
请求持有 R 并等待响应，响应又等待 R，形成 R 到 R 的闭环。把请求与响应放入独立虚拟网络，并规定响应路径不再申请请求网络资源，可拆除该环。
\answercheck{必须保留的检查点}{超时只能检测，不能证明无死锁}
\end{chapteranswer}
